% !TeX root = ../Book1.tex
% 连续同伦的参数方形；每个小格的像落在一个平行标架域中。
\begin{figure}[htbp]
	\centering
	\begin{tikzpicture}[font=\small]
		\fill[GeometryBlueLight] (1.25,1) rectangle (2.5,2);
		\draw[GeometryGray] (0,0) rectangle (5,3);
		\foreach \x in {1.25,2.5,3.75}{
			\draw[GeometryGray] (\x,0)--(\x,3);
		}
		\foreach \y in {1,2}{
			\draw[GeometryGray] (0,\y)--(5,\y);
		}
		\draw[orange!75!black,very thick,->] (1.25,1)--(2.5,1);
		\draw[orange!75!black,very thick,->] (2.5,1)--(2.5,2);
		\draw[GeometryBlue,very thick,->] (1.25,1)--(1.25,2);
		\draw[GeometryBlue,very thick,->] (1.25,2)--(2.5,2);
		\fill[GeometryBlue] (1.25,1) circle[radius=2pt];
		\fill[GeometryBlue] (2.5,2) circle[radius=2pt];
		\node at (2.5,-0.3) {\(H(t,0)=\gamma_0(t)\)};
		\node at (2.5,3.3) {\(H(t,1)=\gamma_1(t)\)};
		\node[left,align=center] at (-0.15,1.5) {左边恒为\\\(p\)};
		\node[right,align=center] at (5.15,1.5) {右边恒为\\\(q\)};
		\draw[->,GeometryGray] (0,-0.85)--(5,-0.85) node[right] {\(t\)};
		\draw[->,GeometryGray] (-2.1,0)--(-2.1,3) node[above] {\(s\)};
	\end{tikzpicture}
	\caption{细分的是同伦的参数方形；网格线的像不必光滑。每个小格的像位于一个平行标架域中，所以蓝色与橙色两条边路产生相同运输。逐格使用这一等式，再利用左右边的运输恒等，便能比较相邻整行。}
	\label{g1:fig:flat-homotopy-grid}
\end{figure}
